\chapter{Introduzione}
Il progetto vuole modellare e implementare un sistema in Java che permetta di generare tracce di un'esecuzione a partire dalla definizione di un taskset.

Ogni traccia è definita come una sequenza di coppie $<tempo, evento>$, dove un $evento$ può essere: rilascio di un job di un task; esecuzione di un job; completamento di un job; preemption che un task può subire.

%%%%%%%%%%%%%%%%%%%%%%%%%%%%%%%%%%%%%%%%%%%%%%%%%%%%%%%%%%%%%%%%
\section{Capacità}
A partire da un taskset, il sistema ha le capacità di:
\begin{itemize}
    \item Generare la traccia di esecuzione del taskset schedulato tramite un dato algoritmo di scheduling (i.e., Rate Monotonic e Earliest Deadline First).
    \item Generare un dataset di tracce relative a più simulazioni su un taskset.
    \item Specificare il tempo desiderato della simulazione.
    \item Rilevare eventuali deadline miss. Se un job non ha finito di eseguire entro la propria deadline allora la simulazione si arresta; non continua perché ci interessa valutare il primo fallimento, visto che i successivi potrebbero essere in cascata del primo.
    \item Campionare da una distribuzione l'execution time di un job.
    \item Definito uno scheduler con il relativo taskset, eseguire il controllo di feasibility adeguato per la schedulabilità del taskset. Per RM è stato implementato l'hyperbolic bound; per EDF il controllo necessario e sufficiente sul fattore di utilizzo.
\end{itemize}